\sectionbreak
\section*{
	\gostTitleFont
	\redline
	ВВЕДЕНИЕ
}
\subtitlespace
{\gostFont
	\par \redline Сервисы заказа такси прочно вошли в повседневную жизнь современного города. Пользователи всё реже ловят автомобиль на улице или звонят диспетчеру, предпочитая оформлять заказ через мобильное приложение. Это требует от серверной части таких сервисов высокой производительности, отказоустойчивости и способности обрабатывать большое количество одновременных запросов.

	\par \redline Данный дипломный проект посвящён разработке серверной части приложения заказа такси. Серверная часть является ключевым компонентом любой подобной платформы: именно на сервере выполняется обработка заказов, поиск свободных водителей, расчёт стоимости поездки, хранение данных о пользователях и координация взаимодействия между всеми участниками процесса.

	\par \redline От качества реализации серверной части напрямую зависит стабильность работы всего приложения. Недостаточная производительность или частые сбои приводят к потере пользователей и снижению конкурентоспособности сервиса. Поэтому проектированию и реализации серверной части должно уделяться особое внимание.

	\par \redline В рамках данной работы разработана серверная часть на основе микросервисной архитектуры. Каждый функциональный блок системы вынесен в отдельный сервис со своей базой данных, а взаимодействие между ними организовано через REST API и брокер сообщений. Такой подход обеспечивает гибкость, независимое развёртывание компонентов и возможность горизонтального масштабирования.

	\par \redline В первом разделе проведён анализ предметной области и сформулированы требования к системе. Во втором разделе спроектирована архитектура системы. Третий раздел описывает реализацию ключевых модулей. Четвертый раздел посвящен тестированию и развёртыванию системы. В пятом разделе выполнен расчёт экономической эффективности.

	\par
}
